\documentclass{standalone}
\usepackage{tikz}
\usetikzlibrary{positioning}

\begin{document}

\begin{tikzpicture}[every node/.style={circle, draw, minimum size=5mm, inner sep=0pt}, node distance=8mm and 12mm]

% 单射
\node[draw=none] at (0,2) {单射 (Injective)};
\node (a1) at (0,1) {A};
\node (a2) at (0,0) {B};
\node (a3) at (0,-1) {C};

\node (b1) at (2,1) {1};
\node (b2) at (2,0) {2};
\node (b3) at (2,-1) {3};

\draw[->] (a1) -- (b1);
\draw[->] (a2) -- (b2);
\draw[->] (a3) -- (b3);

% 满射
\node[draw=none] at (6,2) {满射 (Surjective)};
\node (c1) at (5,1) {A};
\node (c2) at (5,0) {B};

\node (d1) at (7,1.5) {1};
\node (d2) at (7,0.5) {2};
\node (d3) at (7,-0.5) {3};

\draw[->] (c1) -- (d1);
\draw[->] (c1) -- (d2);
\draw[->] (c2) -- (d2);
\draw[->] (c2) -- (d3);

% 双射
\node[draw=none] at (12,2) {双射 (Bijective)};
\node (e1) at (11,1) {A};
\node (e2) at (11,0) {B};
\node (e3) at (11,-1) {C};

\node (f1) at (13,1) {1};
\node (f2) at (13,0) {2};
\node (f3) at (13,-1) {3};

\draw[->] (e1) -- (f2);
\draw[->] (e2) -- (f1);
\draw[->] (e3) -- (f3);

\end{tikzpicture}

\end{document}
